\documentclass[11pt]{article}

\usepackage[margin=1in]{geometry}
\usepackage{parskip}
\usepackage{listings}
\usepackage{xcolor}
\usepackage{booktabs}
\usepackage{hyperref}
\hypersetup{colorlinks=true,linkcolor=blue,urlcolor=blue}

\lstset{
  basicstyle=\ttfamily\small,
  breaklines=true,
  columns=fullflexible,
  frame=single,
  framesep=4pt,
  backgroundcolor=\color{gray!8},
  keywordstyle=\color{blue!70!black}\bfseries,
  commentstyle=\color{green!45!black}\itshape,
  stringstyle=\color{red!60!black},
  language=C,
  showstringspaces=false,
}

\newcommand{\code}[1]{\texttt{\detokenize{#1}}}

\title{Why One \code{swim-times.w} Resolves the Meet and Date and the
Other Does Not}
\author{Comparison of \code{curlswim-save/swim2} and \code{curlswim/swim2}}
\date{\today}

\begin{document}
\maketitle

\section{Summary}

Both programs are the same literate source in every respect except one:
\emph{which USA~Swimming data-hub endpoint they call to retrieve a
swimmer's best times}. That single choice determines whether the swim
date and meet name are available at all.

\begin{itemize}
  \item The version that \textbf{cannot} resolve the meet and date
        (\code{curlswim/swim2}, as committed at \code{HEAD}) calls the
        lightweight \code{GET GetBestTimesForMember} feed. That feed
        returns only stroke, distance, course, and formatted time. It
        carries \emph{no} date, meet, or standard, so the program has
        nothing to parse and deliberately blanks those columns.
  \item The version that \textbf{can} resolve the meet and date
        (\code{curlswim-save/swim2}, and the current working tree of
        \code{curlswim/swim2}) calls the richer
        \code{POST /BestTimes} query --- the same call the data-hub
        front-end issues to render a swimmer's ``Best Times'' page.
        Each returned record includes \code{meetName}, \code{eventCode},
        \code{swimDate}, \code{swimTime}, and \code{timeStandard}, so the
        scanner extracts all five fields.
\end{itemize}

In short: this is not a parsing bug and not a build difference. The
non-working program never had the data in hand because it asked the
wrong endpoint for it.

\section{Verifying the two sources are otherwise identical}

The two \code{.w} files in the working tree are byte-for-byte identical,
as are the tangled \code{.c} and the compiled binaries:

\begin{lstlisting}[language=bash]
$ diff curlswim-save/swim2/swim-times.w curlswim/swim2/swim-times.w
$ md5sum curlswim*/swim2/swim-times.w
f1d02528d7ac764dcc051d42aa93b940  curlswim-save/swim2/swim-times.w
f1d02528d7ac764dcc051d42aa93b940  curlswim/swim2/swim-times.w
\end{lstlisting}

The meaningful difference is between the \emph{committed} version of
\code{curlswim/swim2/swim-times.w} (which used the lightweight feed) and
the fixed version in \code{curlswim-save/swim2} (which the working tree
has since been updated to match). The remainder of this report compares
those two implementations.

\section{The version that fails: \code{GetBestTimesForMember}}

The failing program issues a single \code{GET} per member and reassembles
each event code from the numeric distance, stroke abbreviation, and
course code. Its own prose states plainly that the feed omits the three
fields in question:

\begin{quote}\itshape
``The best-times feed carries only the event and the time, so
\code{date}, \code{meet}, and \code{standard} are left empty --- those
fields are not exposed to anonymous callers.''
\end{quote}

The request and parse loop are:

\begin{lstlisting}
@<Load best-times cache@>=
if (strcmp(g_bt_member, member_id) != 0) {
    free(g_bt_resp);
    char url[512];
    snprintf(url, sizeof url,
             "%s/GetBestTimesForMember/%s", TIMES_API, member_id);
    g_bt_resp = http_request(url, NULL);   /* GET: body is NULL */
    snprintf(g_bt_member, sizeof g_bt_member, "%s", member_id);
}
if (!g_bt_resp) return;

@<Select best time for event@>=
const char *p = g_bt_resp;
char abbr[8], course[8], stime[32];
long dist;
while (nrows < MAX_TIMES &&
       scan_string("strokeAbbreviation", &p, abbr, sizeof abbr)) {
    if (!scan_long("distance", &p, &dist)) break;
    if (!scan_string("courseCode", &p, course, sizeof course)) break;
    if (!scan_string("swimTime", &p, stime, sizeof stime)) break;
    char code[32];
    snprintf(code, sizeof code, "%ld %s %s", dist, abbr, course);
    if (strcmp(code, event_code) != 0) continue;
    rows[nrows].sort_key = time_to_seconds(stime);
    snprintf(rows[nrows].time, sizeof rows[nrows].time, "%s", stime);
    /* No meetName / swimDate / timeStandard in this feed: blank them. */
    rows[nrows].date[0] = rows[nrows].standard[0] = rows[nrows].meet[0] = '\0';
    nrows++;
}
\end{lstlisting}

The decisive line is the explicit
\code{rows[nrows].date[0] = rows[nrows].standard[0] =
rows[nrows].meet[0] = '\textbackslash 0';}. The program scans only
\code{strokeAbbreviation}, \code{distance}, \code{courseCode}, and
\code{swimTime} because those are the only keys the response contains.
There is no \code{meetName} or \code{swimDate} in the JSON to scan for,
so the columns are set empty by construction. No amount of parsing could
recover fields that the server never sent.

\section{The version that works: \code{POST /BestTimes}}

The working program targets the heavier endpoint the front-end itself
uses. Because that query keys on a single stroke-and-distance pair, the
program issues one \code{POST} per distinct pair and caches every parsed
record per member. Crucially, each record now carries the missing
fields:

\begin{lstlisting}
@<Times fetch@>+=
static void fetch_pair(const char *member_id, const char *stroke, long dist)
{
    char url[256], body[256];
    snprintf(url, sizeof url, "%s/BestTimes", TIMES_API);
    snprintf(body, sizeof body,
             "{\"memberId\":\"%s\",\"strokeAbbreviation\":\"%s\","
             "\"distance\":%ld}", member_id, stroke, dist);
    char *resp = http_request(url, body);   /* POST: body is non-NULL */
    if (!resp) return;
    @<Parse BestTimes records@>
    free(resp);
}

@<Parse BestTimes records@>=
const char *p = resp;
char meet[256], event[32], date[16], stime[32], std[48];
while (scan_string("meetName", &p, meet, sizeof meet)) {
    if (!scan_string("eventCode", &p, event, sizeof event)) break;
    if (!scan_string("swimDate",  &p, date,  sizeof date))  break;
    if (!scan_string("swimTime",  &p, stime, sizeof stime)) break;
    std[0] = '\0';
    scan_string("timeStandard", &p, std, sizeof std);
    cache_add(event, stime, date, std, meet);   /* meet + date preserved */
}
\end{lstlisting}

Here the scanner reads \code{meetName} and \code{swimDate} directly out
of each record and hands them to \code{cache_add}, which normalises the
API's \code{"Mon DD, YYYY"} date to ISO \code{YYYY-MM-DD} and stores the
meet name verbatim. The data exists in the response, so it survives into
the output.

\section{Side-by-side}

\begin{center}
\begin{tabular}{@{}p{0.28\textwidth}p{0.32\textwidth}p{0.32\textwidth}@{}}
\toprule
 & \textbf{Fails (\code{curlswim} HEAD)} & \textbf{Works (\code{curlswim-save})} \\
\midrule
Endpoint         & \code{GET GetBestTimesForMember/\{id\}} & \code{POST /BestTimes} \\
HTTP method      & GET (no body)          & POST (JSON body per stroke/distance) \\
Calls per member & One                    & One per distinct stroke-distance pair ($\approx$22) \\
Fields returned  & stroke, distance, course, time & meetName, eventCode, swimDate, swimTime, timeStandard \\
Meet name        & Absent $\Rightarrow$ blanked & Present $\Rightarrow$ parsed \\
Swim date        & Absent $\Rightarrow$ blanked & Present $\Rightarrow$ parsed + ISO-normalised \\
Standard         & Absent $\Rightarrow$ blanked & Present $\Rightarrow$ parsed \\
Helper needed    & \code{scan_long} (numeric distance) & only \code{scan_string} (all values are strings) \\
\bottomrule
\end{tabular}
\end{center}

\section{Root cause}

The failure is an \textbf{endpoint/data-availability} problem, not a
code-quality problem. The \code{GetBestTimesForMember} feed is a
summary view: it returns just enough to reassemble an event code and a
time. The meet name, swim date, and motivational standard are simply not
present in its payload, and the program correctly (if unhelpfully) leaves
those columns empty rather than fabricating them.

The fix, realised in the working version, is to call the richer
\code{POST /BestTimes} query that the data hub's own web front-end uses.
That response includes \code{meetName} and \code{swimDate} (plus
\code{timeStandard}), so the same string scanner --- with no new parsing
machinery beyond dropping the now-unused \code{scan_long} --- can pull
them out. The cost is more network round trips (one POST per
stroke-and-distance pair instead of a single GET), which the program
amortises with a per-member cache.

\end{document}
